% Transcription of Cayley's Collected Mathematical Papers, Volume 4, pages 76-100
% Papers 238 (continuation), 239, 240, 241, 242, 243, 244, 245 (start)

\documentclass[11pt,a4paper]{article}

% ============================================================
% STANDARD ENGLISH MATHEMATICS PREAMBLE
% ============================================================
\usepackage[utf8]{inputenc}
\usepackage[T1]{fontenc}
\usepackage[english,french]{babel}
\usepackage{amsmath,amssymb,amsthm}
\usepackage{mathtools}
\usepackage{geometry}
\usepackage{titlesec}
\usepackage{fancyhdr}
\usepackage{enumitem}
\usepackage{array}
\usepackage{booktabs}
\usepackage[protrusion=true,expansion=false]{microtype}

\geometry{paperwidth=7in,paperheight=10in,left=0.75in,right=0.75in,top=0.85in,bottom=0.85in}

\pagestyle{fancy}
\fancyhf{}
\fancyhead[LE,RO]{\thepage}
\renewcommand{\headrulewidth}{0pt}

\setlength{\abovedisplayskip}{3pt plus 1pt minus 1pt}
\setlength{\belowdisplayskip}{3pt plus 1pt minus 1pt}
\setlength{\abovedisplayshortskip}{0pt plus 1pt}
\setlength{\belowdisplayshortskip}{2pt plus 1pt minus 1pt}

\titleformat{\section}{\normalfont\Large\bfseries}{\thesection.}{0.5em}{}
\titleformat{\subsection}{\normalfont\large\bfseries}{\thesubsection.}{0.5em}{}

\newcommand{\fcn}[2]{(#1\kern-1pt, #2)}
\newcommand{\fxn}[2]{(#1\kern-1pt, #2)}
\DeclareMathOperator{\Pf}{Pf}
\DeclareMathOperator{\sn}{sn}
\DeclareMathOperator{\cn}{cn}
\DeclareMathOperator{\dn}{dn}

\usepackage{graphicx}
\emergencystretch=3em
\tolerance=600
\begin{document}

% ============================================================
% TITLE PAGE
% ============================================================
\thispagestyle{empty}
\begin{center}
\vspace*{2cm}
{\Huge\scshape The Collected Mathematical Papers\par\vspace{0.3em} of\par\vspace{0.5em} Arthur Cayley\par\vspace{1.5em}}
{\LARGE Volume IV\par\vspace{1em}}
{\large Papers 238 (continuation)--245 (part)\par\vspace{0.5em}}
{\large Pages 76--100\par\vspace{2em}}
\vfill
{\large Cambridge University Press\par\vspace{0.3em}}
{\normalsize 1891}
\end{center}
\newpage

% ============================================================
% TABLE OF CONTENTS FOR THIS SECTION
% ============================================================
\tableofcontents
\newpage

\setcounter{page}{76}

% ============================================================
% PAPER 238 (continuation)
% NOTE SUR LES NORMALES D'UNE CONIQUE
% [From Crelle, t. LVI. (1859), pp. 182--185]
% Pages 76--77 of this volume (begins on p. 75, continued here)
% ============================================================

\section*{Note sur les Normales d'une Conique (suite)}
\addcontentsline{toc}{section}{238. Note sur les Normales d'une Conique (suite)}
\label{paper:238cont}

\subsection*{[Continuation from p.~75]}

\selectlanguage{french}

et de ces deux équations on déduit la suivante
\[
\begin{vmatrix}
x_1, & y_1, & z_1 \\
x_2, & y_2, & z_2 \\
x_3, & y_3, & z_3
\end{vmatrix} = 0
\]
en désignant par le symbole qui forme le premier membre la fonction
\[
x_1 y_2 z_3 + x_2 y_3 z_1 + x_3 y_1 z_2 + x_1 y_3 z_2 + x_2 y_1 z_3 + x_3 y_2 z_1.
\]
C'est ce qui résulte de l'identité
\[
\begin{vmatrix}
x_1, & y_1, & z_1 \\
x_2, & y_2, & z_2 \\
x_3, & y_3, & z_3
\end{vmatrix}
\cdot
\begin{vmatrix}
x_1, & y_1, & z_1 \\
x_2, & y_2, & z_2 \\
x_3, & y_3, & z_3
\end{vmatrix}
=
\begin{vmatrix}
x_1^2, & y_1^2, & z_1^2 \\
x_2^2, & y_2^2, & z_2^2 \\
x_3^2, & y_3^2, & z_3^2
\end{vmatrix}
+ 2 x_1 y_1 z_1 \, x_2 y_2 z_2 \, x_3 y_3 z_3
\begin{vmatrix}
\dfrac{1}{x_1}, & \dfrac{1}{y_1}, & \dfrac{1}{z_1} \\[2pt]
\dfrac{1}{x_2}, & \dfrac{1}{y_2}, & \dfrac{1}{z_2} \\[2pt]
\dfrac{1}{x_3}, & \dfrac{1}{y_3}, & \dfrac{1}{z_3}
\end{vmatrix};
\]
car le déterminant qui forme le second facteur du premier membre de l'équation ne s'évanouissant pas, c'est l'autre facteur qui devra s'évanouir en vertu des deux relations données.

L'équation de la droite $(1,2)$ sera
\[
\begin{vmatrix}
x, & y, & z \\
x_1, & y_1, & z_1 \\
x_2, & y_2, & z_2
\end{vmatrix} = 0;
\]
les coordonnées du pôle de cette droite par rapport à la conique donnée $ax^2 + by^2 + cz^2 = 0$, seront
\[
\frac{1}{a}(y_1 z_2 - y_2 z_1) : \frac{1}{b}(z_1 x_2 - z_2 x_1) : \frac{1}{c}(x_1 y_2 - x_2 y_1);
\]
mais les deux équations $a x_1^2 + b y_1^2 + c z_1^2 = 0$, $a x_2^2 + b y_2^2 + c z_2^2 = 0$ donnent $a : b : c = y_1 z_2 - y_2 z_1 : z_1 x_2 - z_2 x_1 : x_1 y_2 - x_2 y_1$; par suite de cela les coordonnées du pôle deviennent
\[
\frac{1}{y_1 z_2 + y_2 z_1} : \frac{1}{z_1 x_2 + z_2 x_1} : \frac{1}{x_1 y_2 + x_2 y_1},
\]
donc l'équation de la polaire (l'harmonicale) de ce pôle par rapport aux trois droites $(x=0,\; y=0,\; z=0)$ sera
\[
x(y_1 z_2 + y_2 z_1) + y(z_1 x_2 + z_2 x_1) + z(x_1 y_2 + x_2 y_1) = 0,
\]
laquelle peut être représentée comme suit
\[
\begin{Bmatrix}
x, & y, & z \\
x_1, & y_1, & z_1 \\
x_2, & y_2, & z_2
\end{Bmatrix} = 0,
\]

\setcounter{page}{77}

et, en substituant $(x_3, y_3, z_3)$ au lieu de $(x, y, z)$, on voit que la droite dont il s'agit contient le point 3. De même cette droite contient le point 4, de sorte que le théorème se trouve démontré.

Observons encore que dans la géométrie de la sphère on peut prendre pour la conique absolue la conique imaginaire qui est l'intersection de la surface sphérique avec la surface conique imaginaire $x^2 + y^2 + z^2 = 0$. Le mot perpendiculaire aura alors la signification ordinaire, et on aura pour les coniques sphériques ce théorème très simple, savoir que les six côtés du quadrilatère inscrit seront les polaires (les harmonicales) --- par rapport aux trois axes de la conique sphérique donnée --- des six sommets du quadrilatère circonscrit. C'est ce que l'on reconnaît aussi par l'analyse que je viens de donner, laquelle en effet est précisément celle dont on se servirait naturellement pour les coniques sphériques.

\medskip
\noindent\textit{Londres, le 10 Mars, 1857.}

\selectlanguage{english}

\newpage

% ============================================================
% PAPER 239
% ADDITION À LA NOTE SUR LA COMPOSITION DU NOMBRE 47
% [From Crelle, t. LVI. (1859), pp. 186--187: Sequel to 236]
% Pages 78--79 of this volume
% ============================================================

\section{Addition à la Note sur la Composition du Nombre 47 par rapport aux Vingt-troisièmes Racines de l'Unité}
\label{paper:239}

\textsc{[From the Journal für die reine und angewandte Mathematik (Crelle), t.~LVI. (1859), pp.~186--187: Sequel to 236.]}

\setcounter{page}{78}

\selectlanguage{french}

M.~Kummer a bien voulu m'écrire une lettre où il remarque que le cube du facteur complexe du nombre 47, multiplié par l'unité complexe convenable, peut effectivement se décomposer en deux facteurs réciproques; c'est-à-dire qu'il existe une unité complexe $E(\alpha)$ telle que
\[
E(\alpha)(\alpha^{11} + \alpha^{10} + \alpha^{8} + \alpha^{7} + \alpha^{4} + \alpha^{2})^{3} = F(\alpha) F(\alpha^{-1}).
\]
Pour $F(\alpha)$ M.~Kummer a trouvé la valeur
\[
F(\alpha) = \alpha^{16} + \alpha^{13} + \alpha^{9} + \alpha^{8} + \alpha^{6} - \alpha^{4} + \alpha^{2},
\]
fonction qui par conséquent est l'un des 22 facteurs complexes de $47^{3}$. Dans un postscriptum M.~Kronecker m'a communiqué l'expression suivante de cette unité complexe
\[
E(\alpha) = \frac{(\alpha^{6} + \alpha^{17})(\alpha^{8} + \alpha^{15})^{2}(\alpha^{10} + \alpha^{13})^{2}(\alpha^{9} + \alpha^{14})}{(\alpha^{5} + \alpha^{18})^{2}}
\]
équivalente à l'expression en fonction entière:
\[
E(\alpha) =
\begin{cases}
-23\alpha + 2\alpha^{2} - 20\alpha^{3} - 21\alpha^{4} - 3\alpha^{5} - 17\alpha^{6} - 4\alpha^{7} - 8\alpha^{8} - 12\alpha^{9} \\
\;\;{}+ \alpha^{10} + \alpha^{11} - 3\alpha^{12} - 3\alpha^{13} - 17\alpha^{14} - 4\alpha^{15} - 14\alpha^{16} - 8\alpha^{17} - 12\alpha^{18} .
\end{cases}
\]

En supposant que cette valeur de $E(\alpha)$ soit connue, on trouve sans peine une condition à laquelle $F(\alpha)$ doit satisfaire. La valeur que j'ai donnée pour $(\alpha^{11} + \alpha^{10} + \alpha^{8} + \alpha^{7} + \alpha^{4} + \alpha^{2})^{2}$

\setcounter{page}{79}

contient le terme constant $+6$; en y ajoutant la quantité évanouissante $-6(1 + \alpha + \ldots + \alpha^{22})$ elle se réduit à
\[
\begin{aligned}
&\;\;\;\;\alpha + \alpha^{2} - 3\alpha^{3} + 3\alpha^{4} + 10\alpha^{5} + 9\alpha^{6} + 3\alpha^{7} + 12\alpha^{8} + 3\alpha^{9} \\
&\;\;\;\;{}+ \alpha^{12} + \alpha^{13} - 3\alpha^{14} + 3\alpha^{15} + 10\alpha^{16} + 9\alpha^{17} + 3\alpha^{18} + 12\alpha^{19} + 3\alpha^{20},
\end{aligned}
\]
et en multipliant cette valeur par $E(\alpha)$, le terme constant sera $+808$; donc en y ajoutant la quantité évanouissante $-808(1+\alpha+\ldots+\alpha^{22})$ on obtient
\[
\begin{aligned}
&E(\alpha)(\alpha^{11} + \alpha^{10} + \alpha^{8} + \alpha^{7} + \alpha^{4} + \alpha^{2})^{3} \\
&= \;\; -5\alpha - 8\alpha^{2} - 7\alpha^{3} - 7\alpha^{4} - 6\alpha^{5} - 3\alpha^{6} - 7\alpha^{7} - 7\alpha^{8} - 5\alpha^{9} \\
&\quad{}- 5\alpha^{10} - 8\alpha^{11} - 7\alpha^{12} - 7\alpha^{13} - 6\alpha^{14} - 3\alpha^{15} - 7\alpha^{16} - 7\alpha^{17} - 5\alpha^{18}.
\end{aligned}
\]
En représentant cette expression par $B\alpha + C\alpha^{2} + \ldots + K\alpha^{18}$, j'écris
\[
B\alpha + C\alpha^{2} + \ldots + K\alpha^{18}
= (\alpha + b\alpha + \ldots + k\alpha^{18})(\alpha + b'\alpha^{-1} + \ldots + k'\alpha^{-18})(1 + \alpha + \ldots + \alpha^{22}),
\]
équation qui subsiste pour $\alpha = 1$. On a donc
\[
B + C + \ldots + K = (a + b + \ldots + k)(a' + b' + \ldots + k') \cdot 23 (a + b + \ldots + k'),
\]
ou, d'après les valeurs de $B$, $C$, \ldots, $K$,
\[
-136 = (a + b + \ldots + k)(a' + b' + \ldots + k'),
\]
ce qui donne
\[
(a + b + \ldots + k)^{2} = -136 \pmod{23}.
\]
On peut ajouter à $(a + b + \ldots + k)^{2}$ un multiple quelconque de $(1 + \alpha + \ldots + \alpha^{22})$ et changer le signe; il est donc permis de prendre $(a + b + \ldots + k)$ positif et plus petit que $\frac{23}{2}$. Cela étant la congruence donne
\[
a + b + \ldots + k = 7.
\]
Au moyen de cette valeur l'équation à laquelle il faut satisfaire devient
\[
\begin{aligned}
7 + 2a &= -a^{2} + 3a^{2} + 4a^{3} - a^{4} + 2a^{5} \\
&\quad{}+ 2a^{6} - a^{7} + 3a^{8} + 4a^{9} - a^{10} + 2a^{11} \\
&= (a + b a \ldots + k a^{18})(a + b a^{-1} \ldots + k a^{-18}).
\end{aligned}
\]
À cause des coefficients numériques négatifs, les coefficients cherchés $a$, $b$, \ldots $k$ ne peuvent pas être tous à la fois positifs; et cela étant il y a quelque manière de satisfaire aux deux conditions ci-dessus écrites, savoir il faut donner à sept des coefficients $a$, $b$, \ldots $k$ les valeurs 1, 1, 1, 1, 1, 1, $-1$, et aux autres coefficients la valeur zéro; l'expression
\[
F(\alpha) = \alpha^{2} + \alpha^{4} + \alpha^{6} + \alpha^{8} + \alpha^{13} - \alpha^{14} + \alpha^{16}
\]
s'accorde en effet avec cette conclusion.

\medskip
\noindent\textit{Londres, le 6 Octobre, 1858.}

\selectlanguage{english}

\newpage

% ============================================================
% PAPER 240
% NOTE ON A THEOREM IN SPHERICAL TRIGONOMETRY
% [From the Philosophical Magazine, vol. XVII. (1859), p. 151]
% Page 80 of this volume
% ============================================================

\section{Note on a Theorem in Spherical Trigonometry}
\label{paper:240}

\textsc{[From the Philosophical Magazine, vol.~XVII. (1859), p.~151.]}

\setcounter{page}{80}

I am not aware that the following theorem has been noticed: viz., in any spherical triangle, if as usual $a$, $b$, $c$ are the sides, and $A$, $B$, $C$ the opposite angles, then
\[
\begin{aligned}
\sin b \sin c + \cos b \cos c \cos A &= \sin B \sin C - \cos B \cos C \cos a, \\
\sin c \sin a + \cos c \cos a \cos B &= \sin C \sin A - \cos C \cos A \cos b, \\
\sin a \sin b + \cos a \cos b \cos C &= \sin A \sin B - \cos A \cos B \cos c.
\end{aligned}
\]
The demonstration is very simple; in fact we have
\[
\begin{aligned}
\sin b \sin c + \cos b \cos c \cos A
 &= \sin b \sin c (\sin^{2} A + \cos^{2} A) + \cos b \cos c \cos A \\
 &= \sin b \sin c \sin^{2} A + \cos A (\cos b \cos c + \sin b \sin c \cos A) \\
 &= \sin B \sin C \sin^{2} a + \cos A \cos a \\
 &= \sin B \sin C (1 - \cos^{2} a) + \cos A \cos a \\
 &= \sin B \sin C + \cos a (\cos A - \sin B \sin C \cos a) \\
 &= \sin B \sin C - \cos B \cos C \cos a,
\end{aligned}
\]
which proves the theorem.

\medskip
\noindent\textit{2, Stone Buildings, W.C., January 5, 1859.}

\medskip
\noindent\footnotesize A geometrical proof and interpretation are given, G.~B.~Airy, ``Remarks on Mr Cayley's Trigonometrical Theorem, etc.,'' \textit{Phil.\ Mag.} same volume, p.~176. I transfer to this place the concluding sentence of the subsequent paper 243. ``I take the opportunity of noticing that the theorem in spherical trigonometry, which I gave in the February Number, is not new, but, as pointed out by Prof.~Chauvenet in the Mathematical Monthly (Cambridge, U.S.), is to be found in Cagnoli's `Trigonometry' (1808).''\normalsize

\newpage

% ============================================================
% PAPER 241
% ON POINSOT'S FOUR NEW REGULAR SOLIDS
% [From the Philosophical Magazine, vol. XVII. (1859), pp. 123--128]
% Pages 81--85 of this volume
% ============================================================

\section{On Poinsot's Four New Regular Solids}
\label{paper:241}

\textsc{[From the Philosophical Magazine, vol.~XVII. (1859), pp.~123--128.]}

\setcounter{page}{81}

It is shown by Poinsot, in the ``Mémoire sur les Polygones et les Polyèdres,'' \textit{Journ.\ École Polyt.}\ vol.~IV.\ pp.~16 to 48 (1810), that, besides the regular polyhedrons of ordinary geometry, there are (of course in an extended signification of the term) four new regular polyhedrons, viz.\ an icosahedron, which I will call the great icosahedron (No.\ 33 of the Memoir), and three dodecahedrons, which I will call the great dodecahedron (No.\ 37), the great stellated dodecahedron (No.\ 38), and the small stellated dodecahedron (No.\ 39). The nature of Poinsot's generalization will be best understood by conceiving, as he does, that the polyhedron is projected on a concentric sphere, so that the faces become spherical polygons. Then for the ordinary polyhedrons of geometry, the sum of the angles at a vertex $= 4$ right angles; but it may, according to the more general notion, be $= s$ times 4 right angles. In like manner for the ordinary polyhedrons, the sides of a face subtend at the centre angles the sum of which is $= e'$ times 4 right angles. And finally, the sum of the spherical polygons is ordinarily equal to the entire spherical surface; but according to the more general notion, it may be $= D$ times the spherical surface ($e$ is Poinsot's $s$; $e'$ does not occur in Poinsot; and, for a reason which will appear, I have written $D$ for Poinsot's $E$.)

The new polyhedra are constructed, as follows:

\subsection*{1.\ The great Icosahedron.}
Each face is made up of seven faces, or rather four faces and six half faces of the ordinary icosahedron, in the manner shown by fig.~1. There are six right angles in each vertex, and together, not four, but eight right angles, or $s=2$; but, as in the ordinary polyhedra, $e'=1$; and the sum of all the faces, or eight times the spherical surface, or $D=7$. (Also $E=7$.)

\setcounter{page}{82}

\subsection*{2.\ The great Dodecahedron.}
Each face is made up of five faces of the ordinary icosahedron in the manner shown by the figure~2. There are five angles at each vertex, and these make up together eight right angles, or $s=2$; but, as in ordinary polyhedra, $e'=1$; and the sum of all the faces is obviously $12 \times \tfrac{5}{12}$, that is three times the spherical surface, or $D=3$. (Also $E=3$.)

\subsection*{3.\ The great stellated Dodecahedron.}
Each face is formed by stellating a face of the great dodecahedron in the manner shown by fig.~3. There are, as in the ordinary

\setcounter{page}{83}

dodecahedron, three angles at each vertex, and the sum of these is simply four right angles, or $s=1$. On account of the stellation, $e'=2$. Each of the projecting parts of the face is equal $\tfrac{1}{6}$ of the face of the ordinary icosahedron; and if we reckon the area of the stellated pentagon as twice the interior pentagon plus the projecting parts, then the area of the stellated face will be four times the spherical surface, and accordingly Poinsot writes $E=4$. If, however, what seems preferable, we reckon the area of the stellated pentagon as a side (or what is the same thing, reckon the stellated face as the centre of the face and standing on the sides.

\subsection*{4.\ The small stellated Dodecahedron.}
Each face is twelve that of the face of the ordinary dodecahedron, and the sum of the faces is twice the spherical surface, and accordingly Poinsot writes $E = 2$. But according to the second mode of measurement, the area of the stellated face is three times the spherical surface; and the sum of the faces is three times the spherical surface, and the sum is $D=3$.

\setcounter{page}{84}

$E$, viz.\ the faces make together $E$ times the spherical surface, the area of a stellated face being reckoned as the sum of the triangles having their vertices at the centre of the face and standing on the sides.

$D$, viz.\ the faces make together $D$ times the spherical surface, the area of a stellated face reckoned as the sum of the triangles having their vertices at the centre of the face and standing on the sides.

The Table is

\begin{center}
\begin{tabular}{lcccccccc}
\toprule
Designation & $H$ & $S$ & $A$ & $a$ & $a'$ & $e$ & $e'$ & $D$ \\
\midrule
Tetrahedron      & 4  & 4  & 6  & 3 & 3 & 1 & 1 & 1 \\
Hexahedron       & 6  & 8  & 12 & 4 & 3 & 1 & 1 & 1 \\
Octahedron       & 8  & 6  & 12 & 3 & 4 & 1 & 1 & 1 \\
Dodecahedron     & 12 & 20 & 30 & 5 & 3 & 1 & 1 & 1 \\
Icosahedron      & 20 & 12 & 30 & 3 & 5 & 1 & 1 & 1 \\
\midrule
Great stellated dodecahedron     & 12 & 20 & 30 & 5 & 3 & 1 & 2 & 7 \\
Great icosahedron                & 20 & 12 & 30 & 3 & 5 & 2 & 1 & 7 \\
Small stellated dodecahedron     & 12 & 12 & 30 & 5 & 5 & 1 & 2 & 3 \\
Great dodecahedron               & 12 & 12 & 30 & 5 & 5 & 2 & 1 & 3 \\
\bottomrule
\end{tabular}
\end{center}

where the figures which are polar reciprocals of each other are written in pairs: viz., as is well known, the tetrahedron is its own reciprocal, the hexahedron and octahedron are reciprocals, and the dodecahedron and icosahedron are reciprocals; moreover the great stellated dodecahedron and the great dodecahedron are reciprocals. I have not been able to define $E$ in such a manner as to enable me to form the definition of a reciprocal number $E'$; this may be possible, but in the mean time it seems better to discard $E$ altogether, and use instead of it the number $D$.

Euler's well-known relation applying to ordinary polyhedra is
\[
S + H = A + 2.
\]
Poinsot in his memoir has (by an extension of Legendre's demonstration of Euler's theorem) obtained the more general relation,
\[
eS + H = A + 2E,
\]

\setcounter{page}{85}

which, however, does not apply to the two stellated figures where $e'$ is different from unity; the general form is
\[
eS + e'H = A + 2D,
\]
which applies to all the nine figures. This applies to all polyhedra, regular or not, which are such that $e$ has the same value for each vertex, and $e'$ the same value for each face. If for any other polyhedra, regular or not, $e$, then, according to the foregoing mode of reckoning, the area of the face (measured in the ordinary manner) is to be reckoned as $e'$ times the same value $a$, then, according to the foregoing mode of reckoning, the area of the face will be $E + e'H = A + 2D$.

I remark that the small stellated dodecahedron and the great dodecahedron are connected by the following Table:

\begin{center}
\begin{tabular}{lll}
$a\;b\;c\;d\;e=P$, & $A\;C\;E\;R\;D=p$, \\
$p\;b\;i\;k\;e=A$, & $P\;I\;E\;B\;H=a$, \\
$p\;e\;j\;i\;a=B$, & $P\;J\;A\;C\;I=b$, \\
$p\;d\;f\;j\;b=C$, & $P\;F\;B\;D\;J=c$, \\
$p\;c\;g\;f\;c=D$, & $P\;G\;C\;E\;F=d$, \\
$p\;a\;h\;g\;d=E$, & $P\;H\;D\;A\;G=e$, \\
$j\;c\;d\;g\;q=F$, & $J\;D\;Q\;C\;G=f$, \\
$f\;d\;e\;h\;q=G$, & $F\;E\;Q\;D\;H=g$, \\
$g\;e\;a\;i\;q=H$, & $G\;A\;Q\;E\;I=h$, \\
$h\;a\;b\;j\;q=I$, & $H\;B\;Q\;A\;J=i$, \\
$i\;b\;c\;f\;q=J$, & $I\;C\;Q\;B\;F=j$, \\
$f\;g\;h\;i\;j=Q$, & $F\;H\;J\;G\;I=q$,
\end{tabular}
\end{center}
where it is to be noticed that in either Table each non-consecutive triad of any printed column occurs once, and only once, as a non-consecutive triad of another printed. The restriction that a non-consecutive triad of any multiplet is not to occur as a triad, consecutive or non-consecutive, of any other multiplet (see my note appended to Mr Kirkman's paper ``On Autopolar Polyhedra,'' \textit{Phil.\ Trans.}\ 1857, p.~183 [259]), applies only to ordinary polyhedra, and not to the class here considered.

\medskip
\noindent\textit{2, Stone Buildings, W.C., January 13, 1859.}

\newpage

% ============================================================
% PAPER 242
% SECOND NOTE ON POINSOT'S FOUR NEW REGULAR POLYHEDRA
% [From the Philosophical Magazine, vol. XVII. (1859), pp. 209--210]
% Pages 86--87 of this volume
% ============================================================

\section{Second Note on Poinsot's Four New Regular Polyhedra}
\label{paper:242}

\textsc{[From the Philosophical Magazine, vol.~XVII. (1859), pp.~209--210.]}

\setcounter{page}{86}

The Note on Poinsot's four new regular Polyhedra (Feb.\ Number, p.~123), [241] was written without my being acquainted with Cauchy's first memoir, ``Recherches sur les Polyèdres'' (\textit{Journ.\ Polyt.}\ vol.~IX.\ pp.~68--86, 1813), the former part of which (pp.~68--76) relates to Poinsot's polyhedra. Cauchy considers the polyhedra, not as projected on the sphere, but as solids, and he shows, very elegantly, that all such polyhedra are obtained by stellating the regular dodecahedron and the regular icosahedron. The reciprocal method would lead to two figures, and explained and developed in my former Note. Cauchy does not at all consider Poinsot's generalized equation, $eS + H = A + 2E$; but the latter part of his memoir employed by Poinsot, and explained and developed in my former Note. Cauchy does not at all consider Poinsot's generalized equation, $eS + H = A + 2E$; viz.\ Cauchy's theorem is, for any polyhedron, regular or not, partitioned into any number of polyhedra (including those of the original polyhedron), and $A$ the total number of edges (including those of the original polyhedron), and $A$ the total number of vertices (including those of the original polyhedron), then $S + H = A + P + 1$; that is, the sum of the number of vertices and the number of faces exceeds by unity the sum of the number of edges and the number of polyhedra.

For $P=1$, we have Euler's equation $S + H = A + 2$; and we have a theorem relating to the partition of a polygon; viz.\ if the polygon is divided into $R$ polygons, and if $S$ be the number of vertices, and $A$ the number of edges, then $S + H = A + 1$;

\setcounter{page}{87}

polyhedra. I remark that, in the equation $S + H = A + 1$, $H$ should, in analogy with Cauchy's notation for polyhedra, be replaced by $P$; so that we have for a single polygon,
\[
A = S;
\]
and for the partitions of a polygon,
\[
A = S + P - 1;
\]
corresponding respectively to Euler's theorem for a single polyhedron, viz.
\[
S + H = A + 2;
\]
and to Cauchy's theorem for the partitions of a polyhedron, viz.
\[
S + H = A + 2 + (P - 1).
\]
Cauchy's second memoir (pp.~87--98) contains a very beautiful demonstration of the theorem implied in the ninth definition of the eleventh book of Euclid, viz.\ that two convex polyhedra are equal when they are bounded by the same number of faces equal each to each.

\medskip
\noindent\textit{2, Stone Buildings, W.C., February 1, 1859.}

\newpage

% ============================================================
% PAPER 243
% ON THE THEORY OF GROUPS, AS DEPENDING ON THE
% SYMBOLIC EQUATION $\theta^n = 1$. THIRD PART
% [From the Philosophical Magazine, vol. XVIII. (1859), pp. 34--37: Sequel to 125 and 126]
% Pages 88--91 of this volume
% ============================================================

\section{On the Theory of Groups, as Depending on the Symbolic Equation $\theta^{n} = 1$. Third Part}
\label{paper:243}

\textsc{[From the Philosophical Magazine, vol.~XVIII. (1859), pp.~34--37: Sequel to 125 and 126.]}

\setcounter{page}{88}

The following is, I believe, a complete enumeration of the groups of 8:
\begin{center}
\begin{tabular}{rl}
I.   & $1$, $\alpha$, $\alpha^{2}$, $\alpha^{3}$, $\alpha^{4}$, $\alpha^{5}$, $\alpha^{6}$, $\alpha^{7}$ \quad ($\alpha^{8} = 1$). \\
II.  & $1$, $\alpha$, $\alpha^{2}$, $\alpha^{3}$, $\beta$, $\beta\alpha$, $\beta\alpha^{2}$, $\beta\alpha^{3}$ \quad ($\alpha^{4} = 1$, $\beta^{2} = 1$, $\alpha\beta = \beta\alpha$). \\
III. & $1$, $\alpha$, $\alpha^{2}$, $\alpha^{3}$, $\beta$, $\beta\alpha$, $\beta\alpha^{2}$, $\beta\alpha^{3}$ \quad ($\alpha^{4} = 1$, $\beta^{2} = 1$, $\alpha\beta = \beta\alpha^{3}$). \\
IV.  & $1$, $\alpha$, $\alpha^{2}$, $\alpha^{3}$, $\beta$, $\beta\alpha$, $\beta\alpha^{2}$, $\beta\alpha^{3}$ \quad ($\alpha^{4} = 1$, $\beta^{2} = \alpha^{2}$, $\alpha\beta = \beta\alpha^{3}$). \\
V.   & $1$, $\alpha$, $\beta$, $\beta\alpha$, $\gamma$, $\gamma\alpha$, $\gamma\beta$, $\gamma\beta\alpha$ \quad ($\alpha^{2} = 1$, $\beta^{2} = 1$, $\gamma^{2} = 1$, $\alpha\beta = \beta\alpha$, $\alpha\gamma = \gamma\alpha$, $\beta\gamma = \gamma\beta$).
\end{tabular}
\end{center}

That the groups are really distinct is perhaps most readily seen by writing down the indices of the different terms of each group; these are
\begin{center}
\begin{tabular}{rl}
I.   & $1, 8, 4, 8, 2, 8, 4, 8$. \\
II.  & $1, 4, 2, 4, 2, 4, 2, 4$. \\
III. & $1, 4, 2, 4, 2, 2, 2, 2$. \\
IV.  & $1, 4, 2, 4, 4, 4, 4, 4$. \\
V.   & $1, 2, 2, 2, 2, 2, 2, 2$.
\end{tabular}
\end{center}

It will be presently seen why there is no group where the symbols $\alpha$, $\beta$ are such that $\alpha^{4} = 1$, $\beta^{2} = 1$, $\alpha\beta = \beta\alpha^{2}$. A group which presents itself for consideration is
\[
1,\; \alpha,\; \alpha^{2},\; \alpha^{3},\; \beta,\; \beta\alpha,\; \beta\alpha^{2},\; \beta\alpha^{3} \quad (\alpha^{4} = 1,\; \beta^{2} = \alpha^{2},\; \alpha\beta = \beta\alpha);
\]

\setcounter{page}{89}

but the indices of the different terms of this group are
\[
1,\; 4,\; 2,\; 4,\; 2,\; 4,\; 2,\; 4,
\]
and if we write $\beta\alpha = \gamma$, then we find $\gamma^{2} = \beta\alpha\beta\alpha = \beta\beta\alpha\alpha = \alpha^{4} = 1$, $\alpha\gamma = \alpha\beta\alpha = \beta\alpha\alpha = \gamma\alpha$; and the group is
\[
1,\; \alpha,\; \alpha^{2},\; \alpha^{3},\; \gamma,\; \gamma\alpha,\; \gamma\alpha^{2},\; \gamma\alpha^{3} \quad (\alpha^{4} = 1,\; \gamma^{2} = 1,\; \alpha\gamma = \gamma\alpha),
\]
which is the group II.

The group IV is a remarkable one; it appears to arise from the circumstance that the factors 2 and 4 of the number 8 are not prime to each other; this can only happen when the number which denotes the order of the group contains a square factor. But the nature of the group in question will be better understood by presenting it under a different form. In fact, if we write $\beta\alpha^{3} = \gamma$, $\alpha^{2} = \beta^{2} = \delta$, then we find $\alpha = \delta\alpha$, $\beta\alpha^{2} = \delta\beta$, $\beta\alpha = \delta\gamma$, and the group will be
\[
1,\; \alpha,\; \beta,\; \gamma,\; \delta,\; \delta\alpha,\; \delta\beta,\; \delta\gamma,
\]
where the laws of combination are
\[
\begin{aligned}
&\delta^{2} = 1,\quad \alpha^{2} = \beta^{2} = \gamma^{2} = \delta, \quad \alpha\beta = \gamma, \\
&\beta\gamma = \alpha,\quad \gamma\alpha = \beta, \quad \beta\alpha = \gamma\delta = \delta\gamma, \\
&\gamma\beta = \alpha\delta = \delta\alpha, \quad \alpha\gamma = \beta\delta = \delta\beta.
\end{aligned}
\]

Observe that $\delta$ is a symbol of operation such that $\delta^{2} = 1$, and that $\delta$ is convertible with each of the other symbols $\alpha$, $\beta$, $\gamma$. It will be not so much a restrictive assumption in regard to $\delta$, as a definition of $-1$ considered as a symbol of operation if we write $\delta = -1$; the group thus becomes
\[
1,\; \alpha,\; \beta,\; \gamma,\; -1,\; -\alpha,\; -\beta,\; -\gamma,
\]
where
\[
\alpha = \beta\gamma = -\gamma\beta, \quad \beta = \gamma\alpha = -\alpha\gamma, \quad \gamma = \alpha\beta = -\beta\alpha.
\]
Hence $\alpha$, $\beta$, $\gamma$ combine according to the laws of the quaternion symbols $i$, $j$, $k$; and it is only the point of view from which the question is here considered which obliges us to consider the symbols as belonging to a group of 8, instead of (as in the theory of quaternions) a group of 4.

Suppose in general that the symbols $\alpha$, $\beta$ are such that
\[
\alpha^{m} = 1,\quad \beta^{n} = 1,\quad \alpha\beta = \beta\alpha^{s},
\]
then we find
\[
\alpha^{u}\beta^{v} = \beta^{v}\alpha^{u s^{v}};
\]

\setcounter{page}{90}

and therefore if $v = n$, $\alpha^{u} = \alpha^{u s^{n}}$ or $\alpha^{u(s^{n} - 1)} = 1$, whence $u(s^{n} - 1) \equiv 0 \pmod{m}$; or since $u$ is arbitrary, $s^{n} - 1 \equiv 0 \pmod{m}$, an equation which, if $m$, $n$ are given, determines the admissible values of $s$; thus, for example, if $n = 2$, and $m$ is a prime number, then $s = 1$ or $s = m - 1$. The equation $\alpha^{u}\beta^{v} = \beta^{v}\alpha^{u s^{v}}$ shows that any combination whatever of the symbols $\alpha$, $\beta$ can be expressed in the form $\beta^{q}\alpha^{p}$ (or, if we please, in the form $\alpha^{p}\beta^{q}$). It is proper to show that the assumed law is consistent with the associative law, viz.\ that the expression
\[
\beta^{v_{1}}\alpha^{u_{1}} \cdot \beta^{v_{2}}\alpha^{u_{2}} \cdot \beta^{v_{3}}\alpha^{u_{3}}
\]
can be transformed in one way only into the form $\beta^{v}\alpha^{u}$. We in fact have
\[
\beta^{v_{1}}\alpha^{u_{1}} \cdot \beta^{v_{2}}\alpha^{u_{2}} = \beta^{v_{1}+v_{2}}\alpha^{u_{1}s^{v_{2}}+u_{2}},
\]
and multiplying this by the remaining factor $\beta^{v_{3}}\alpha^{u_{3}}$, we have
\[
\beta^{v_{1}+v_{2}+v_{3}}\alpha^{(u_{1}s^{v_{2}}+u_{2})s^{v_{3}} + u_{3}},
\]
which is equal to
\[
\beta^{v_{1}+v_{2}+v_{3}}\alpha^{u_{1}s^{v_{2}+v_{3}}+u_{2}s^{v_{3}}+u_{3}};
\]
and the result would have been precisely the same if, instead of thus combining together the first and second factors and the product with the third factor, we had combined the first factor with the product of the second and third factors, so that the associative law is satisfied.

It is now easy to see that if, as before,
\[
\alpha^{m} = 1,\quad \beta^{n} = 1,\quad \alpha\beta = \beta\alpha^{s},
\]
conditions which it has been shown imply $s^{n} \equiv 1 \pmod{m}$, then the symbols $\beta^{q}\alpha^{p}$ (or, if we please, $\alpha^{p}\beta^{q}$), where $p$ has the values $0, 1, 2, \ldots, m-1$, and $q$ the values $0, 1, 2, \ldots, n-1$, form a group of $nm$ terms. In particular, as already noticed, if $n=2$ and $m$ is prime, then $s = 1$ or $s = m-1$; the two groups so obtained are essentially distinct from each other. If $n = 2$, but $m$ is not prime, then $s$ has in general more than two values: thus for $m=12$, $s^{2} \equiv 1 \pmod{12}$, which is satisfied by $s = 1, 5, 7$ and $11$; the group corresponding to $s = 1$ is distinct from that for any other value of $s$, but I have not ascertained whether the values other than unity do, or do not, give groups distinct from each other.

For the sake of an observation to which it gives rise, I write down an example of a group corresponding to $n = 2$, $s = m-1$, say $m = 5$, and therefore $s = 4$, so that we have
\[
\alpha^{5} = 1,\quad \beta^{2} = 1,\quad \alpha\beta = \beta\alpha^{4},
\]
and the group is
\[
1,\; \alpha,\; \alpha^{2},\; \alpha^{3},\; \alpha^{4},\; \beta,\; \beta\alpha,\; \beta\alpha^{2},\; \beta\alpha^{3},\; \beta\alpha^{4},
\]

\setcounter{page}{91}

the indices of the several terms being
\[
1,\; 5,\; 5,\; 5,\; 5,\; 2,\; 2,\; 2,\; 2,\; 2.
\]

The group is here expressed by means of the symbols $\alpha$, $\beta$, having the indices 5 and 2 respectively, but it may be expressed by means of two symbols having each of them the index 2. Thus putting $\beta\alpha = \gamma$, we find $\beta^{2} = 1$, $\gamma^{2} = 1$, $(\beta\gamma)^{5} = 1$, which is equivalent to $(\gamma\beta)^{5} = 1$, and the group may be represented in the form
\[
1,\; \beta,\; \gamma,\; \beta\gamma,\; \gamma\beta,\; \beta\gamma\beta,\; \gamma\beta\gamma,\; \beta\gamma\beta\gamma,\; \gamma\beta\gamma\beta,\; \beta\gamma\beta\gamma\beta = \gamma\beta\gamma\beta\gamma,
\]
the equality of the last two symbols being an obvious consequence of the equation $(\beta\gamma)^{5} = 1$. It is clear that for any even number $2p$ whatever, there is always a group which can be expressed in this form.

\medskip
\noindent\textit{2, Stone Buildings, W.C., June 9, 1859.}

\newpage

% ============================================================
% PAPER 244
% ON AN ANALYTICAL THEOREM RELATING TO THE DISTRIBUTION
% OF ELECTRICITY UPON SPHERICAL SURFACES
% [From the Philosophical Magazine, vol. XVIII. (1859), pp. 119--127]
% Pages 92--98 of this volume
% ============================================================

\section{On an Analytical Theorem relating to the Distribution of Electricity upon Spherical Surfaces}
\label{paper:244}

\textsc{[From the Philosophical Magazine, vol.~XVIII. (1859), pp.~119--127.]}

\setcounter{page}{92}

There is contained in Plana's ``Mémoire sur la distribution de l'électricité à la surface de deux sphères conductrices complètement isolées'' \textit{(Mém.\ de Turin}, vol.~VII.\ 1845), an identical relation which is remarkable, as well in itself as because by means of it the author corrects an error into which Poisson had fallen in his researches on the same subject. The development of a certain definite integral is obtained in the form (equation 165)
\[
y = -2 \cdot 3 \frac{M_{1} h}{b^{2}} \sin^{4}\tfrac{1}{2}\theta + \frac{3 \cdot 4 \cdot 5}{1 \cdot 2}\frac{M_{2} h}{b^{3}} \sin^{6}\tfrac{1}{2}\theta + \&\mathrm{c}.
\]
Poisson had in effect shown that $M_{1} = 0$; and he thence inferred that, $\theta$ being small, the function in question $\propto \sin^{4}\tfrac{1}{2}\theta$, or what is the same thing, $\propto (1 - \cos\theta)^{2}$. In the former part of the memoir, Plana shows that this is not the true form of the development; the foregoing development must therefore be illusory; and Plana in fact shows, by a laborious induction carried as far as $M_{7}$, that all the coefficients $M$ vanish identically. The identical equation $M_{i} = 0$, where $i$ is any positive integer whatever, constitutes the analytical theorem above referred to. Plana's expression for the function $M_{i}$ is as follows:

$B_{1}$, $B_{2}$, $B_{3}$, \&c.\ denote Bernoulli's numbers as given by the equation
\[
\frac{t}{e^{t} - 1} = 1 - \tfrac{1}{2} t - B_{1}\frac{t^{2}}{1 \cdot 2} + B_{2}\frac{t^{4}}{1 \cdot 2 \cdot 3 \cdot 4} + \&\mathrm{c}.
\]
($B_{1} = \tfrac{1}{6}$, $B_{2} = \tfrac{1}{30}$, $B_{3} = \tfrac{1}{42}$, $B_{4} = \tfrac{1}{30}$, \&c. I have, in conformity with the usual practice, written the equations so as to make these numbers all positive; with Plana they are

\setcounter{page}{93}

alternately positive and negative). And in the equation 162, writing for $k$ its value $\frac{1}{1+b}$, we have, $\lambda$ being any positive integer,
\[
\begin{aligned}
G_{\lambda} &= \frac{1}{\lambda+1}-\tfrac{1}{2}(1+b) + B_{1}\frac{\lambda}{1 \cdot 2}\frac{(1+b)^{2}}{\lambda - 1} \\
&\quad{} - B_{2}\frac{\lambda \cdot \lambda - 1 \cdot \lambda - 2}{1 \cdot 2 \cdot 3 \cdot 4} \frac{(1+b)^{4}}{\lambda - 3} \\
&\quad{} + B_{3}\frac{\lambda \cdot \lambda - 1 \cdot \lambda - 2 \cdot \lambda - 3 \cdot \lambda - 4}{1 \cdot 2 \cdot 3 \cdot 4 \cdot 5 \cdot 6}\frac{(1+b)^{6}}{\lambda - 5} \\
&\quad{} + \&\mathrm{c}.,
\end{aligned}
\]
where the series is continued for so long as the factor in the denominator is positive. It should be observed that this factor really divides out, and that the rule just mentioned amounts to this, viz.\ that when $\lambda$ is odd, the finite series on the right-hand side is to be continued to its last term; but when $\lambda$ is even, the series is to be continued only to the last term but one. And $G_{\lambda}$ being thus defined, the expression for $M_{i}$ (see equation 164, in which I have written for $k$ its value $\frac{1}{1+b}$) is
\[
M_{i} = (1+b)^{i-1}\left\{ G_{i} + (1+2)\frac{i \cdot i-1}{1 \cdot 2} \left(\frac{1}{b}\right)^{2} G_{i+2} + (1+3)\frac{i \cdot i-1 \cdot i-2 \cdot i-3}{1 \cdot 2 \cdot 3 \cdot 4} \left(\frac{1}{b}\right)^{4} G_{i+4} + \&\mathrm{c}.\right\}
\]
where on the right-hand side the finite series is continued up to its last term, the value of which is obviously
\[
\frac{1}{i}(1+2i-1)\left(\frac{1}{b}\right)^{2i-1} G_{3i-1},\quad \text{that is},\quad \frac{2}{i}\left(\frac{1}{b}\right)^{2i-1} G_{3i-1}.
\]
But the form of this equation may be somewhat simplified. We in fact have
\[
G_{i} = \frac{1}{i+1} - \tfrac{1}{2}(1+b) + B_{1}\frac{i}{1 \cdot 2}\frac{(1+b)^{2}}{i-1} - B_{2}\frac{i \cdot i-1 \cdot i-2}{1 \cdot 2 \cdot 3 \cdot 4}\frac{(1+b)^{4}}{i-3} + \&\mathrm{c.},
\]
which is at once changed into
\[
(1+b) M_{i} = \;\theta_{i} + \tfrac{i+1}{i}\,\tfrac{i}{1 \cdot 2}\,\theta_{i+1} + \&\mathrm{c.},
\]
where
\[
\theta_{i} = \tfrac{1}{i+1}-\tfrac{1}{2}i(1+b) + B_{1}\frac{i \cdot i-1}{1 \cdot 2}(1+b)^{2} - B_{2}\frac{i \cdot i-1 \cdot i-2 \cdot i-3}{1 \cdot 2 \cdot 3 \cdot 4}(1+b)^{4} + \&\mathrm{c.}
\]

\setcounter{page}{94}

or multiplying by $i + 1$, and then putting $i - 1$ in the place of $i$, we have
\[
i(1+b)M_{i-1} = \theta_{i} + \frac{i \cdot 1}{1 \cdot 2}\theta_{i+1} + \frac{i \cdot i-1}{1 \cdot 2 \cdot 3}\theta_{i+2} + \&\mathrm{c.} = 0,
\]
where
\[
\theta_{i} = 1 - \tfrac{1}{2}i(1+b) + B_{1}\frac{i \cdot i-1}{1 \cdot 2}(1+b)^{2} - B_{2}\frac{i \cdot i-1 \cdot i-2 \cdot i-3}{1 \cdot 2 \cdot 3 \cdot 4}(1+b)^{4} - \&\mathrm{c.}
\]
In this last equation the finite series on the right-hand side is, when $i$ is even, to be continued up to its last term, but when $i$ is odd, then only up to the last term but one. The equation to be proved is
\[
0 = \theta_{i} + \frac{i \cdot 1}{1 \cdot 2}\theta_{i+1} + \frac{i \cdot i-1}{1 \cdot 2 \cdot 3}\theta_{i+2} + \&\mathrm{c.},
\]
where on the right-hand side the finite series is to be continued up to its last term: and the equation holds for any integer value of $i$ which is $\geq 2$. This is the simplest form of Plana's theorem.

We have
\[
\theta_{i} = i(1+b)^{i-1} G_{i-1};
\]
or writing this equation under the form $\theta_{i} = i(1+b)\frac{G_{i-1}}{(1+b)^{2-i}}$, and comparing with Plana's developed expressions for $G$ (which are continued by him as far as $G_{11}$), we find
\[
\begin{aligned}
\theta_{2} &= -\tfrac{1}{6}\,b, \\
\theta_{3} &= \tfrac{1}{6}\,b^{2}, \\
\theta_{4} &= \tfrac{1}{6}\,b + \tfrac{1}{2}b^{2} - \tfrac{1}{6}b^{3} - \tfrac{1}{6}b^{4}, \\
\theta_{5} &= -\tfrac{2}{15}\,b^{2} - 2 b^{3} - \tfrac{5}{6}b^{5}, \\
\theta_{6} &= -\tfrac{1}{6}\,b - \tfrac{1}{2}b^{2} - \tfrac{2}{3}b^{3} + \tfrac{4}{3}b^{4} + b^{5} + \tfrac{1}{6}b^{6}, \\
\theta_{7} &= \tfrac{2}{15}\,b^{2} + 4 b^{3} + \tfrac{15}{2}b^{4} + 4 b^{5} + \tfrac{2}{15}b^{6}, \\
\theta_{8} &= \tfrac{2}{15}\,b + \tfrac{1}{2}b^{2} + \tfrac{5}{3}b^{3} + \tfrac{i}{6}b^{4} + \tfrac{15}{6}b^{5} + b^{6} + \tfrac{1}{4}b^{7} + \tfrac{1}{30}b^{8}, \\
&\&\mathrm{c.} = -\tfrac{2}{15}b^{2} - 12 b^{3} - 37 b^{4} - 54 b^{5} - 37 b^{6} - 12 b^{7} - \tfrac{2}{15}b^{8},
\end{aligned}
\]

\setcounter{page}{95}

which are of course the results obtained by developing the foregoing expression for $\theta_{i}$, in powers of $b$, and collecting the terms. The formulae put in evidence a remarkable symmetry which does not exist in the original expression in powers of $1+b$.

It would be now easy to verify, for moderately small values of the suffix, the equations
\[
\theta_{2} + \tfrac{2}{1 \cdot 2}\theta_{3} + \tfrac{2 \cdot 1}{1 \cdot 2 \cdot 3}\theta_{4} = 0,
\]
\[
\&\mathrm{c.}
\]
This is, in fact, Plana's process, which, however, as the suffixes increase, becomes a very laborious one, and the law of the terms which destroy each other is not in anywise exhibited thereby.

I have succeeded in obtaining a complete demonstration, founded on Herschel's theorem for the development of a function of $e^{t}$, and the expression thereby given for Bernoulli's numbers. The theorem in question [See Herschel's \textit{Collection of Examples in the Calculus of Finite Differences}, Cambridge, 1820, p.~70, where the theorem is given in the form $f\{(1 + \Delta)^{x}\}\, 0^{x} = u^{x} f(1 + \Delta)\, 0^{x}$] is, that for any function of $e^{t}$ which admits of development in positive integer powers of $t$,
\[
f(e^{t}) = f(1 + \Delta) e^{t \cdot 0},
\]
where the right-hand side denotes the series the general term whereof is
\[
\frac{t^{n}}{1 \cdot 2 \cdot 3 \ldots n}\, f(1 + \Delta)\, 0^{n};
\]
and $f(1 + \Delta)$ is of course to be developed in powers of $\Delta$, $\Delta^{2}$, $\Delta^{3}$, \&c., applied to the symbol $0^{n}$ (viz.\ $\Delta 0^{n} = 1^{n} - 0^{n}$; $\Delta^{2} 0^{n} = 2^{n} - 2 \cdot 1^{n} + 0^{n}$, \&c.).

This gives
\[
\frac{t}{e^{t} - 1} = \frac{\log(1 + \Delta)}{\Delta}\, e^{t \cdot 0},
\]
and comparing the development of the right-hand side with the development
\[
\frac{t}{e^{t} - 1} = 1 - \tfrac{1}{2}t + B_{1}\frac{t^{2}}{1 \cdot 2} - B_{2}\frac{t^{4}}{1 \cdot 2 \cdot 3 \cdot 4} + \&\mathrm{c.},
\]
we find
\[
\frac{\log(1 + \Delta)}{\Delta}\, 0^{0} = 1,
\]
\[
\frac{\log(1 + \Delta)}{\Delta}\, 0 = -\tfrac{1}{2},
\]
\[
\frac{\log(1 + \Delta)}{\Delta}\, 0^{2\sigma} = (-)^{\sigma - 1} B_{\sigma},
\]
\[
\frac{\log(1 + \Delta)}{\Delta}\, 0^{2\sigma + 1} = 0, \quad (\sigma \geq 1).
\]

\setcounter{page}{96}

It is now easy to obtain the equation
\[
\theta_{i} = \frac{\log(1 + \Delta)}{\Delta}\,\{(1 + 0(1+b))^{i} - (-0(1+b))^{i}\}.
\]
In fact, the first two terms of the development of the expression on the right-hand side agree with those of the foregoing expression for $\theta_{i}$. For any even power $2x$ (except, when $i$ is even, the power $2x = i$) the term is
\[
\frac{[i]^{2x}}{[2x]^{2x}}(1 + b)^{2x}\frac{\log(1 + \Delta)}{\Delta}\, 0^{2x},
\]
which agrees; and when $i$ is even, then for the power $2x = i$ there are two equal and opposite terms which destroy each other, and the whole term in $\theta_{i}$ is, as it ought to be, zero. For any odd power $2x - 1$, $(x > 1)$, (including, when $i$ is odd, the power $2x - 1 = i$), the term vanishes as containing an evanescent factor. The expression for $\theta_{i}$ is thus shown to be true.

I write for shortness,
\[
X = 1 + 0(1 + b),\quad Y = -0(1 + b).
\]
Forming the expression for $i(1+b) M_{i-1}$,
\[
\theta_{i} + \frac{i \cdot 1}{1 \cdot 2}\theta_{i+1} + \frac{i \cdot i+1}{1 \cdot 2 \cdot 3}\theta_{i+2} + \&\mathrm{c.},
\]
this is
\[
i(1+b)M_{i-1} = \frac{\log(1+\Delta)}{\Delta}\left\{\left(X\!\left(1+\tfrac{X}{b}\right)\right)^{i} - \left(Y\!\left(1+\tfrac{Y}{b}\right)\right)^{i}\right\},
\]
and we have
\[
X = (1 + 0)(1 + b) - b,
\]
and therefore
\[
1 + \frac{X}{b} = \frac{1+0}{b}\,(1+b),
\]
\[
\left(X\left(1+\tfrac{X}{b}\right)\right)^{i} = \frac{(1+0)^{i}(1+b)^{i}}{b^{i}}\,(1+0)\,\{(1+0)(1+b) - b\}^{i},
\]
\[
1 + \frac{Y}{b} = -\frac{0(1+b)-b}{b} = -\frac{1}{b}\{0(1+b) - b\}.
\]

\setcounter{page}{97}

We see that the expression for $\left(X\left(1+\tfrac{X}{b}\right)\right)^{i}$ is deduced from that of $\left(Y\left(1+\tfrac{Y}{b}\right)\right)^{i}$ by writing therein $1 + 0$ in the place of $0$; we have therefore
\[
\left(X\left(1+\tfrac{X}{b}\right)\right)^{i} = \Delta_{0}\left(Y\left(1+\tfrac{Y}{b}\right)\right)^{i}
\]
and consequently
\[
\left(X\left(1+\tfrac{X}{b}\right)\right)^{i} - \left(Y\left(1+\tfrac{Y}{b}\right)\right)^{i} = (\Delta_{0} - 1)\left(Y\left(1+\tfrac{Y}{b}\right)\right)^{i} = \log(1+\Delta_{0})\, (0\{0(1+b)-b\})^{i};
\]
whence also
\[
i(1+b) M_{i-1} = \left(\frac{1}{b}\right)^{i}\,\log(1+\Delta)\,(0\{0(1+b) - b\})^{i}.
\]
We have by the general theorem,
\[
t = \log e^{t} = \log(1 + \Delta) e^{t \cdot 0};
\]
and consequently whenever $n \geq 2$,
\[
\log(1 + \Delta)\, 0^{n} = 0.
\]
But if $i \geq 2$, the reduced function $(0\{0(1+b) - b\})^{i}$ contains only powers of $0$, and the superior powers; it is therefore reduced to zero by the operation $\log(1+\Delta)$, and we have
\[
i(1+b) M_{i-1} = \theta_{i} + \frac{i \cdot 1}{1 \cdot 2}\theta_{i+1} + \frac{i \cdot i-1}{1 \cdot 2 \cdot 3}\theta_{i+2} + \&\mathrm{c.} = 0;
\]
and the theorem in question is thus proved. The foregoing expressions for $\theta_{2}$, $\theta_{3}$, \&c.\ show that these functions all divide by $b$, and moreover that when $i$ is even and greater than 2, then that $\theta_{i}$ divides by $b^{2}$. The equation
\[
\theta_{i} = \frac{\log(1 + \Delta)}{\Delta}\,\{(1 + 0(1+b))^{i} - (-0(1+b))^{i}\}
\]
gives generally for the term in $\theta_{i}$ involving $b^{\alpha}$, the expression
\[
\frac{[i]^{\alpha}}{[\alpha]^{\alpha}}\,\frac{\log(1+\Delta)}{\Delta}\,\{(1+0)^{i-\alpha} 0^{\alpha} - (-0)^{i}\};
\]
and it is to be shown, first, that the coefficient vanishes for $\alpha = 0$; and next, that when $i$ is even and $> 2$, the coefficient also vanishes for $\alpha = 1$. Putting $\alpha = 0$, the coefficient is
\[
\frac{\log(1+\Delta)}{\Delta}\,\{(1+0)^{i} - (-0)^{i}\},
\]
or to
\[
\log(1+\Delta)\, 0^{i-1} + \{1 - (-)^{i}\}\frac{\log(1+\Delta)}{\Delta}\, 0^{i},
\]

\setcounter{page}{98}

where, since $i \geq 2$, the former term vanishes, as above remarked; and the latter term, when $i$ is even, vanishes on account of the factor $1 - (-)^{i}$; and when $i$ is odd, on account of the other factor. Hence the coefficient vanishes for $\alpha = 0$.

Next, if $i$ is even, and $\alpha = 1$, the coefficient becomes
\[
\frac{\log(1+\Delta)}{\Delta}\,\{(1+0)^{i-1} 0 - (-0)^{i}\},
\]
which, writing $(1 + 0) - 1$ for $0$, becomes
\[
\frac{\log(1+\Delta)}{\Delta}\,\{(1 + 0)^{i} - (1 + 0)^{i-1} - 0^{i}\},
\]
which, since $(1+0)^{i} - 0^{i} = \Delta 0^{i}$, $(1+0)^{i-1} = (1+\Delta)\, 0^{i-1}$, is equal to
\[
\log(1+\Delta)\, 0^{i-1} - \frac{(1+\Delta) \log(1+\Delta)}{\Delta}\, 0^{i-1},
\]
or since the first term vanishes, to
\[
-\frac{(1+\Delta) \log(1+\Delta)}{\Delta}\, 0^{i-1}.
\]
But this function is to a numerical factor $\frac{1}{i}$ près the coefficient of $t^{i-1}$ in $\frac{e^{t}}{e^{t} - 1}$, or (what is the same thing) in $\frac{t}{e^{-t} - 1}$; and if in the expression for $\frac{t}{e^{t} - 1}$ we write $-t$ in the place of $t$, we find
\[
\frac{t}{e^{-t} - 1} = -1 - \tfrac{1}{2}t + B_{1}\frac{t^{2}}{1 \cdot 2} - B_{2}\frac{t^{4}}{1 \cdot 2 \cdot 3 \cdot 4} + \&\mathrm{c.},
\]
so that, $i$ being even and greater than 2, the function in question vanishes. Hence in the case under consideration the coefficient vanishes for $\alpha = 1$.

Writing $\beta$ for $i - \alpha$, or assuming $\alpha + \beta = i$, the symmetry of the foregoing expressions for $\theta_{2}$, $\theta_{3}$, \&c.\ shows that we ought to have
\[
\frac{\log(1+\Delta)}{\Delta}\,\{(1+0)^{\alpha} 0^{\beta} - (-0)^{\alpha + \beta}\} = \pm \frac{\log(1+\Delta)}{\Delta}\,\{(1+0)^{\beta} 0^{\alpha} - (-0)^{\alpha + \beta}\},
\]
where the upper or under sign is to be taken according as $\alpha + \beta$ is even or odd. Or separating the two cases, we find
\[
\frac{\log(1+\Delta)}{\Delta}\,\{(1+0)^{\alpha} 0^{\beta} - (1+0)^{\beta} 0^{\alpha}\} = 0,\quad \alpha+\beta\ \text{even},
\]
\[
\frac{\log(1+\Delta)}{\Delta}\,\{(1+0)^{\alpha} 0^{\beta} + (1+0)^{\beta} 0^{\alpha} - 2 \cdot 0^{\alpha + \beta}\} = 0,\quad \alpha+\beta\ \text{odd}.
\]
I have not attempted to verify \textit{à posteriori} these elegant formulae.

\medskip
\noindent\textit{2, Stone Buildings, W.C., June 18, 1859.}

\newpage

% ============================================================
% PAPER 245
% ON AN ANALYTICAL THEOREM CONNECTED WITH THE DISTRIBUTION
% OF ELECTRICITY ON SPHERICAL SURFACES. SECOND PART
% [From the Philosophical Magazine, vol. XVIII. (1859), pp. 193--202: continuation of 244]
% Pages 99--100 of this volume (continues to p. 107)
% ============================================================

\section{On an Analytical Theorem connected with the Distribution of Electricity on Spherical Surfaces. Second Part}
\label{paper:245}

\textsc{[From the Philosophical Magazine, vol.~XVIII. (1859), pp.~193--202: continuation of 244.]}

\setcounter{page}{99}

The theorem is certainly true; but its existence gives rise to a difficulty to which I shall advert in the sequel. I propose, in the first instance, to give a demonstration which starts from the expression for $fx$ given by Plana's equation (115), instead of the deduced equation which was the basis of my former proof. It will be proper to explain the origin and meaning of the formulae. We have two conducting spherical surfaces, radii $1$ and $b$, in contact with each other (so that the distance between the centres is $1 + b$): and then, if $x$ is the distance from the centre of the sphere, radius $1$, of an exterior point, and $\mu\;(= \cos\theta)$ the cosine of the inclination of this distance to the line from the centre to the centre of the other sphere, the potential $\phi(\mu, x)$ of the sphere, radius $1$, at the point whose coordinates are $(x, \mu)$ is deduced from the potential $fx$ of a point in the axis; that is, if
\[
fx = A_{0} + A_{1}\, x + A_{2}\, x^{2} + \&\mathrm{c.},
\]
then
\[
\phi(\mu, x) = A_{0} P_{0} + A_{1} P_{1}\, x + A_{2} P_{2}\, x^{2} + \&\mathrm{c.},
\]
where $P_{0}$, $P_{1}$, $P_{2}$, \&c.\ are Legendre's functions, viz.\ the functions of $\mu$ which are the coefficients of the successive powers of $x$ in the development of $(1 - 2\mu x + x^{2})^{-\frac{1}{2}}$ in ascending powers of $x$. And the electrical thickness $y$ at any point of the surface of the sphere, radius $1$, is given by the formula
\[
y = x\,\frac{d\phi(\mu, x)}{dx} + \tfrac{1}{2}\phi(\mu, x),
\]
where, after the differentiation, $x = 1$.

\setcounter{page}{100}

The problem consequently depends on the determination of the potential $fx$ for a point on the axis; and this is determined by the functional equation
\[
fx = \frac{h}{1 + 2b - (1+b)x}\, F\!\left(\frac{1 + b - x}{1 + 2b - (1+b)x}\right) + \frac{hb}{1+b-x},
\]
(Plana's equation (G), in which I have written for $\beta$, $\gamma$, $H$ their values, and substituted also for $g$ its value $= h$). The solution of this equation is (equation (H), writing therein $g = h$)
\[
fx = \sum_{0}^{\infty}\frac{h}{n + n(1+b) - n(1+b)x} - \sum_{0}^{\infty}\frac{Ph}{(n+1)(1+b) - (1 + n(1+b))x},
\]
where $P$ is an arbitrary constant \textit{quoad} the functional equation, viz.\ it is any function whatever which has the property of remaining unaltered when $x$ is changed into $\frac{1+b-x}{1+2b-(1+b)x}$. Poisson, and Plana after him, arrive at the conclusion that in the physical problem $P = 0$. It appears to me that there is ground for holding that this is only true \textit{sub modo}, and that $\frac{P}{1+b-x}$ for $x = 1$ (which, if $P$ were retained, would be a term occurring in the expression for the thickness at the point of contact) is not of necessity zero. But the term, if it exists, can be replaced at the conclusion; and I write therefore
\[
fx = \sum_{0}^{\infty}\frac{h}{n + n(1+b) - n(1+b)x} - \sum_{0}^{\infty}\frac{Ph}{(n+1)(1+b) - (1 + n(1+b))x}.
\]
According to the process by which the solution of the functional equation was obtained, this is the true form of the solution; for although the series are non-convergent, and the two sums are in fact each of them infinite, there is nothing to show a relation between the number of terms which must be taken in each series. However, nothing immediately turns upon this, as the expression is only used for obtaining an expression for $fx$ in the form of a definite integral, viz., equation (36),
\[
fx = \frac{b}{1 + 2b - (1+b)x}\int_{0}^{1}\frac{dt(t^{-\frac{b}{1+b}} - 1)\, t^{\frac{bx}{(1+b)(1-x)}}}{1 - t}\,,
\]
or, equation (39),
\[
fx = \frac{b}{1+b}\int_{0}^{1}\frac{1}{1 - t(1+b)(1-x)}\,.
\]

\end{document}
